\documentclass[11pt]{article}
\usepackage{amsmath,amssymb,amsthm,hyperref}
\newtheorem{theorem}{Theorem}
\newtheorem{lemma}[theorem]{Lemma}
\newtheorem{proposition}[theorem]{Proposition}
\newtheorem{corollary}[theorem]{Corollary}
\title{Nonideal Root Variations: Compilation, Oscillation, and Low-Level Fibers}
\author{}
\date{November 23, 2026}
\begin{document}
\maketitle

\paragraph{Scope.}
Nonconstant dependence on the auxiliary variable can escape the common
Laplace-sign obstruction. This note gives an exact compiler and an
example of that escape. A different low-level-fiber invariant still
excludes fixed-root constructions of sublinear auxiliary degree.
These are manual proofs. The universal head bound remains unproved.

\section{An explicit compiler without an ideal-difference hypothesis}

Let $U_h,V_h,a_h,b_h$ be affine forms, $1\le h\le m$.
Assume that all coefficients of every $U_h,V_h$ are strictly positive,
and that at every Boolean vertex the two base lists $U,V$ have equal
multisets of values. The root values may vary with the input. Define
\[
 L(x)=c-\sum_h\frac{a_h(x)}{U_h(x)^2}
          +\sum_h\frac{b_h(x)}{V_h(x)^2}.
 \tag{1}
\]

\begin{proposition}[Finite positive-slope realization]
Suppose $y(x)L(x)\ge\gamma>0$ at every cube vertex. Then, for every
sufficiently small positive rational $\epsilon$, the score
\[
 S_\epsilon(x)=c+\frac1\epsilon
       \left(\sum_h\frac1{U_h(x)+\epsilon a_h(x)}
                  -\sum_h\frac1{V_h(x)+\epsilon b_h(x)}\right)
 \tag{2}
\]
strictly represents $y$ using $2m$ admissible heads. More explicitly,
it suffices to make every perturbed coefficient positive and to have
\[
 \max_x\sum_h\left(
 \frac{\epsilon a_h(x)^2}{U_h(x)^2(U_h(x)+\epsilon a_h(x))}
 +\frac{\epsilon b_h(x)^2}{V_h(x)^2(V_h(x)+\epsilon b_h(x))}
                  \right)<\gamma.
 \tag{3}
\]
\end{proposition}
\begin{proof}
For $D,D+\epsilon a>0$, direct multiplication gives
\[
 \frac1\epsilon\left(\frac1{D+\epsilon a}-\frac1D\right)
       =-\frac a{D^2}
                    +\frac{\epsilon a^2}{D^2(D+\epsilon a)}.
 \tag{4}
\]
The base reciprocal sums in (2) cancel exactly, by the multiset
hypothesis. Apply (4) to each summand. The triangle inequality bounds
$|S_\epsilon-L|$ by the expression in (3). It tends uniformly to zero
on the finite cube. Positivity of the finitely many base coefficients
persists for sufficiently small $\epsilon$. Choosing a rational
$\epsilon$ in that interval proves all assertions. The numerators in
(2) are constants $\epsilon^{-1}$ and $-\epsilon^{-1}$, and the free
bias is allowed by the normal form.
\end{proof}

The script \texttt{search\_nonideal\_root\_variation.py} implements
(3) using exact rational arithmetic after a candidate is found. It
halves $\epsilon$ until the perturbed coefficients are positive and
the bound is less than half the leading margin. An independent checker
then evaluates every actual fraction at every Boolean vertex.
Solver timeouts or failures are not lower bounds.

\begin{proposition}[A nonideal variation with no common aligned time]
Put $B(x)=2+2x$ and $d(x)=1-2x$ on the one-bit cube. The score
\[
 L(x)=d(x)\left(\frac1{B(x)^2}
                          -\frac2{(B(x)+1)^2}\right)
 \tag{5}
\]
is positive at both inputs, but its Laplace integrands have no common
strictly positive time. It is a special case of (1), with $m=2$.
\end{proposition}
\begin{proof}
The two values are $L(0)=1/36$ and $L(1)=7/400$.
Its integrand is
\[
 t d(x)e^{-tB(x)}(1-2e^{-t}).
\]
At $x=0$ this is positive exactly when $t>\log 2$; at $x=1$ it is
positive exactly when $0<t<\log 2$. At $t=\log 2$ both vanish.
Take $U=V=(B,B+1)$, $a=(-d,2d)$, and $b=(0,0)$ in (1).
The common base product is $P(s)=(s+B)(s+B+1)$, and its first
product difference is $d(x)(s+B(x)-1)$, which has degree one in $s$.
Thus the example escapes precisely the ideal, degree-zero hypothesis.
It is a boundary example, not an efficient universal representation.
\end{proof}

\section{Auxiliary degree controls Laplace oscillation}

\begin{lemma}[Zeros of exponential polynomials]
Let $b_1,\ldots,b_r$ be distinct real numbers and let $p_j$ be real
polynomials of degree at most $m_j-1$, with $m_j\ge1$ and
$M=\sum_jm_j$. A nonzero function
$g(t)=\sum_jp_j(t)e^{-b_jt}$ has at most $M-1$ real zeros, counted
with multiplicity, on every finite interval.
\end{lemma}
\begin{proof}
Induct on $M$. Multiply by $e^{b_1t}$, which does not change zeros,
and differentiate. The block at $b_1$ loses one dimension, while
each other polynomial block retains its degree bound. The derivative
therefore belongs to a space of the stated form with total dimension
$M-1$. If it is nonzero, Rolle's theorem with multiplicities and the
induction hypothesis give the result. If it is zero, the original
function is a constant multiple of $e^{-b_1t}$ and has no zeros.
The case $M=1$ is immediate. Zero polynomial blocks can be removed
whenever they arise in the induction.
\end{proof}

\begin{theorem}[A numerator-degree oscillation bound]
Let $P(s)=\prod_{h=1}^m(s+b_h)$ with every $b_h>0$, permitting
repetitions. Let $R$ be nonzero of degree $d<m$. The inverse Laplace
transform $K_R$ of $R/P$, in the class of exponential polynomials,
has at most $d$ zeros on $(0,\infty)$, counted with multiplicity.
Consequently the inverse transform of $(R/P)'$ has at most $d$
positive zeros as well.
\end{theorem}
\begin{proof}
Partial fractions express $K_R$ as an exponential polynomial of total
dimension at most $m$. If the leading coefficient of $R$ is $r_d$,
expansion at infinity gives
\[
       \frac{R(s)}{P(s)}=r_ds^{-(m-d)}+O(s^{-(m-d+1)}).
\]
Repeated integration by parts in the Laplace transform gives
\[
 K_R^{(j)}(0)=0\quad(0\le j<m-d-1),\qquad
 K_R^{(m-d-1)}(0)=r_d\ne0.
\]
For completeness, each remainder after integration by parts is
$O(s^{-k-1})$ because every derivative of the finite exponential
polynomial is bounded near zero and decays at infinity. Comparing
successive powers of $1/s$ justifies the displayed equalities.
The zero at zero has multiplicity $m-d-1$. The preceding lemma leaves
at most $d$ zeros on the positive axis. If there were more, a finite
interval containing zero and those zeros would contradict that lemma.
Finally the transform of $-tK_R(t)$ is $(R/P)'$, by differentiating
under the absolutely convergent integral. Multiplication by $-t$
does not change the positive zeros.
\end{proof}

If $F_\epsilon-G_\epsilon=\epsilon^kR+O(\epsilon^{k+1})$ and
$F_0=G_0=P$, then
\[
 \frac{(F_\epsilon)_s}{F_\epsilon}
       -\frac{(G_\epsilon)_s}{G_\epsilon}
       =\epsilon^k(R/P)'+O(\epsilon^{k+1}).
 \tag{6}
\]
Thus at each input, a leading product difference of auxiliary degree
at most $d$ has a leading exponential integrand with at most $d$
positive zeros. This bound does not assert a common aligned time:
(5) already shows why that assertion fails for $d=1$.

\section{Low root levels have short coordinate paths}

\begin{lemma}[Low-level fibers of an affine form]
Suppose a real affine form $u$ takes values in a fixed ordered set
$\lambda_1<\cdots<\lambda_m$ on the Boolean cube. Complement its
negative-slope coordinates and discard its zero-slope coordinates,
so that
\[
              u(x)=\lambda_r+\sum_{i\in I}w_i z_i,
              \qquad w_i>0.
\]
If $u(x)=\lambda_j$, then at most $j-r$ of the active coordinates
$z_i$ equal one. In particular the first $q$ level fibers can be
nonempty only within $q-1$ coordinate steps of the minimum face.
\end{lemma}
\begin{proof}
If $k$ active coordinates equal one, turn on those coordinates one at
a time from a point of the minimum face. The $k+1$ successive values
are strictly increasing members of the ordered root set, beginning
at $\lambda_r$ and ending at $\lambda_j$. Hence $k+1\le j-r+1$.
\end{proof}

The interval-cardinality version of this argument is kernel checked in
\texttt{Results/LowRootFibers.lean}, as
\texttt{HeadComplexity.card\_add\_one\_le\_affineLevelInterval\_card}
and \texttt{HeadComplexity.activeCoordinates\_lt\_prefixRootCount}.
The counting arguments below are manual, not Lean declarations.

\begin{lemma}[Counting all first-level-fiber patterns]
For $1\le q\le m$, put
\[
 M_q=q\sum_{i=0}^{\min(q-1,n)}\binom ni,
 \qquad
 T_q=(q+1)3^n(1+2M_q)^{n+q+1}.
 \tag{7}
\]
As both the affine form and the ordered root set vary, there are at
most $T_q$ different ordered lists of its fibers at
$\lambda_1,\ldots,\lambda_q$, provided its entire cube image lies
in that root set.
\end{lemma}
\begin{proof}
The signs and zero status of the slopes have at most $3^n$ choices.
There are $q$ possible minimum indices $r\le q$ and one further
case in which all the first $q$ fibers are empty. By the preceding
lemma, only active-coordinate subsets of size at most $q-1$ need
be considered. For these subsets and each $j\le q$, the equalities
\[
                   a+\sum_{i\in S}w_i-\lambda_j=0
\]
are $M_q$ linear tests in at most $n+q+1$ real parameters.
Their complete sign patterns, including zeros, determine all the
fibers; ignoring positivity and all other root constraints only
enlarges this count.

An arrangement of $M$ affine hyperplanes in dimension $k$ has at
most $\sum_{i=0}^k2^i\binom Mi\le(1+2M)^k$ sign patterns when
zeros are retained. To see the first bound, add a hyperplane. Each
face of the induced arrangement on it contributes at most one new
face there and one additional split face, giving the recurrence
$F(M,k)\le F(M-1,k)+2F(M-1,k-1)$. The binomial expression satisfies
that recurrence and the boundary cases. This proves (7).
\end{proof}

\section{Sublinear auxiliary degree is insufficient with fixed roots}

Let each of two affine lists $u_1,\ldots,u_m$ and $v_1,\ldots,v_m$
permute the same fixed distinct roots $\lambda_1<\cdots<\lambda_m$
at every input. The list and roots may depend on the target. Let
$a_h,b_h$ be arbitrary affine perturbations, and suppose
\[
 R(s,x)=\left.\partial_\epsilon\right|_0
   \left[\prod_h(s+u_h(x)+\epsilon a_h(x))
         -\prod_h(s+v_h(x)+\epsilon b_h(x))\right]
 \tag{8}
\]
has degree at most $q-1$ in $s$, at every cube input, where
$1\le q\le m$. Write $P(s)=\prod_{j=1}^m(s+\lambda_j)$.
Allow an arbitrary affine common shift $C(x)$ with
$C(x)+\lambda_j>0$, and an arbitrary bias $c$. Consider the strict
signs of
\[
                   c+\left.(R/P)'\right|_{s=C(x)}.
 \tag{9}
\]
An actual head realization additionally requires the slope-sign
condition for $C+u_h$ and $C+v_h$. We omit that restriction when
counting, thereby enlarging the family. No change of $C$ after
fixing a sign table is assumed to preserve that table.

\begin{theorem}[Fixed-root low-degree sign count]
Put $N=2^n$, $k_1=2m(n+1)+1$, $k_2=m+n+1$, and $p=k_1+k_2$.
When $p\le N$, the number of strict sign tables in (9) is at most
\[
 T_q^{2m}\left(\frac{32768N}{p}\right)^p
                  2^{k_1}(2m)^{k_2}.
 \tag{10}
\]
For any fixed $C_0>0$ and any positive integer sequence $q(n)=o(n)$, the union
over $m\le C_0 2^n/n^2$ and auxiliary degrees at most $q(n)-1$
has only $2^{o(2^n)}$ strict sign tables. It therefore cannot be a
universal construction at that head budget.
\end{theorem}
\begin{proof}
At a vertex let
\[
 D_j(x)=\sum_{h:u_h(x)=\lambda_j}a_h(x)
                   -\sum_{h:v_h(x)=\lambda_j}b_h(x).
\]
Evaluation of (8) at $s=-\lambda_j$ gives
$R(-\lambda_j,x)=D_j(x)\prod_{\ell\ne j}(\lambda_\ell-\lambda_j)$.
Interpolation at the first $q$ roots therefore yields the polynomial
identity in $s$
\[
 R(s,x)=\sum_{j=1}^q D_j(x)
       \prod_{\ell=q+1}^m(\lambda_\ell-\lambda_j)
       \prod_{\substack{1\le\ell\le q\\\ell\ne j}}(s+\lambda_\ell).
 \tag{11}
\]
The interpolation denominators have canceled exactly. No rational
dependence on the root parameters is left in (11).

The first $q$ fibers of the $2m$ affine forms have at most $T_q^{2m}$
possibilities. Fix them. Then every $D_j(x)$ is linear in the
$2m(n+1)$ numerator parameters. After multiplying (9) by the positive
$P(C(x))^2$, its numerator is
\[
       cP(C(x))^2+R'(C(x),x)P(C(x))-R(C(x),x)P'(C(x)).
\]
For each input this is a polynomial, of degree at most one in the
$k_1$ numerator-and-bias parameters and at most $2m$ in the $k_2$
root-and-shift parameters. Apply the two-block polynomial sign bound
already stated and sourced in \texttt{mechanism-audit.tex}, the
largest-term and shifted-Sidon counting arguments. Enlarging the
parameter domain by dropping the root ordering and admissibility
restrictions only increases the count. This proves (10).

For $q=o(n)$, eventually $q\le n/2$. The elementary binomial bound
gives
\[
 \log M_q=O\bigl(q\log(en/q)+\log n\bigr)=o(n),
 \qquad \log T_q=o(n^2).
\]
This includes $q=1$ directly. Thus $2m\log T_q=o(N)$ at the
stated budget. The other terms in the logarithm of (10) are
$O_{C_0}(N\log n/n+n^2)=o(N)$, as in the earlier two-block count.
For sufficiently large $n$, $p\le N$ throughout the budget.
Taking the union over at most $C_0N/n^2$ possible $m$ adds only
$O_{C_0}(n)$ to the logarithm. If $m<q(n)$, replace $q(n)$ by $m$;
this does not worsen any bound.
\end{proof}

\begin{corollary}[Analytic parameter order does not remove this restriction]
The same fixed-root count applies if (8) is replaced by the first
nonzero product difference at any finite order of analytic affine
root perturbations, provided the base roots are distinct and fixed
across the cube and its auxiliary degree is at most $q-1$.
\end{corollary}
\begin{proof}
Suppose $F_\epsilon-G_\epsilon=\epsilon^kR+O(\epsilon^{k+1})$.
Fix an input and match base roots, writing $u_h(0)=v_{\pi(h)}(0)$.
Evaluate the product difference at $s=-u_h(\epsilon)$. The resulting
value of $G_\epsilon$ is $O(\epsilon^k)$. All its factors except
$v_{\pi(h)}(\epsilon)-u_h(\epsilon)$ stay bounded away from zero,
because the base roots are distinct. Thus the matched roots have
equal Taylor coefficients through order $k-1$.

Let $a_h,b_h$ be their coefficients at order $k$. They are affine in
the input: the constant and linear coefficients of an affine form
can be recovered from its values at zero and at the unit vectors,
so their Taylor coefficients are affine as well. The coefficient at
order $k$ of the product difference is exactly (8), since all
lower matched root coefficients agree. The preceding theorem applies.
This argument concerns a common finite order $k$; it does not assert
that every arbitrary head representation has such a description.
\end{proof}

\paragraph{Boundary of the result.}
The theorem fixes the base root multiset across inputs; it does not
assume that the roots are fixed across targets. Moving root multisets,
auxiliary degree comparable to $n$ or larger, and general admissible
head scores remain outside this low-fiber count. No universal existence
lemma for those remaining classes is asserted. The side computation
examines moving-root nonideal variations and is not evidence for an
asymptotic theorem.

\end{document}
